\documentclass[11pt]{article}
\usepackage[a4paper,margin=22mm,headheight=15pt]{geometry}
\usepackage[T1]{fontenc}
\usepackage[utf8]{inputenc}
\usepackage{lmodern}
\usepackage{microtype}
\usepackage{amsmath,amssymb,mathtools,bm}
\usepackage{booktabs,longtable,array,tabularx}
\usepackage{enumitem}
\usepackage{xcolor}
\usepackage{hyperref}
\usepackage{fancyhdr}
\usepackage{listings}
\usepackage{caption}
\usepackage{float}
\usepackage{needspace}
\usepackage{siunitx}

\definecolor{navy}{RGB}{25,54,93}
\definecolor{bluegray}{RGB}{71,93,116}
\definecolor{lightblue}{RGB}{235,242,249}
\definecolor{lightgreen}{RGB}{236,247,239}
\definecolor{lightamber}{RGB}{255,247,226}
\definecolor{lightred}{RGB}{252,237,237}
\definecolor{darkgreen}{RGB}{28,105,63}
\definecolor{darkamber}{RGB}{150,98,0}
\definecolor{darkred}{RGB}{155,42,42}

\hypersetup{colorlinks=true,linkcolor=navy,citecolor=navy,urlcolor=bluegray,
  pdftitle={Analytic Stress in DFT+eDMFT with LAPW: Implemented Equations and Derivation},
  pdfauthor={Technical working note prepared for the eDMFT project}}
\setlength{\parindent}{0pt}
\setlength{\parskip}{5pt plus 1pt}
\setlist{nosep,leftmargin=1.6em}
\allowdisplaybreaks
\numberwithin{equation}{section}

\pagestyle{fancy}
\fancyhf{}
\fancyhead[L]{\small\color{bluegray}Analytic stress in DFT+eDMFT with LAPW}
\fancyhead[R]{\small\color{bluegray}Implemented equations and current audit}
\fancyfoot[C]{\thepage}
\renewcommand{\headrulewidth}{0.3pt}

\newcommand{\Tr}{\operatorname{Tr}}
\newcommand{\Ree}{\operatorname{Re}}
\newcommand{\diag}{\operatorname{diag}}
\newcommand{\sgn}{\operatorname{sgn}}
\newcommand{\dd}{\mathrm d}
\newcommand{\ii}{\mathrm i}
\newcommand{\kb}{\mathbf k}
\newcommand{\one}{\mathbb I}
\newcommand{\eps}{\varepsilon}
\newcommand{\Sig}{\widetilde\Sigma}
\newcommand{\norm}[1]{\left\lVert #1\right\rVert}
\newcommand{\code}[1]{\texttt{#1}}
\newcommand{\filepath}[1]{\nolinkurl{#1}}
\newcommand{\status}[2]{\par\medskip\noindent\colorbox{#1}{\parbox{0.965\linewidth}{#2}}\par\medskip}
\newcommand{\implemented}[1]{\status{lightgreen}{\textbf{Implemented.} #1}}
\newcommand{\proposed}[1]{\status{lightamber}{\textbf{Proposed diagnostic / not yet implemented.} #1}}
\newcommand{\warningbox}[1]{\status{lightred}{\textbf{Important qualification.} #1}}
\newcommand{\identitybox}[1]{\status{lightblue}{\textbf{Stationary identity.} #1}}

\lstset{basicstyle=\ttfamily\small,frame=single,backgroundcolor=\color{lightblue},
  rulecolor=\color{bluegray},breaklines=true,columns=fullflexible}

\begin{document}

\begin{titlepage}
\centering
\vspace*{22mm}
{\Huge\bfseries\color{navy}Analytic Stress in DFT+eDMFT with LAPW\par}
\vspace{7mm}
{\LARGE Implemented Equations, Bird's-Eye Derivation,\par
and the Direct-$C_{\kb}$ Diagnostic\par}
\vspace{13mm}
{\large A self-contained technical working note\par}
\vspace{18mm}
\begin{minipage}{0.86\textwidth}
\small
This note records the equations presently implemented in \code{dmft2}, their
relationship to the stationary DFT+eDMFT Klein functional and to the WIEN2k
LAPW stress kernels, and the remaining diagnostic needed when the projector
metric is singular at individual $\kb$ points.  Every symbol is defined, and
implemented equations are separated from exact target identities and proposed
tests.  The emphasis is on the generalized off-diagonal kinetic,
incomplete-basis, and overlap matrices required by an interacting density
matrix.
\end{minipage}
\vfill
{\large 5 November 2026\par}
\vspace{7mm}
{\small Based on the DFT+eDMFT force formalism, the Blaha--Belbase--Tr\"oster
LAPW stress derivation, and the current \code{eDMFT\_new/src/dmft2} implementation.\par}
\end{titlepage}

\tableofcontents
\clearpage

\section{Purpose, scope, and status labels}

The goal is to compute the first derivative of the stationary DFT+eDMFT
thermodynamic functional under a homogeneous strain.  The derivation must not
differentiate the nonstationary total-energy expression.  It must instead
differentiate the Klein functional at its stationary point.  This removes the
implicit response of $G$ and $\Sigma$---and hence the two-particle vertex---but
retains every \emph{explicit} dependence on the LAPW basis, its metric, and the
normalized correlated projector.

Four kinds of statement are distinguished throughout:

\begin{longtable}{>{\bfseries}p{0.20\linewidth}p{0.72\linewidth}}
\toprule
Stationary identity & An exact one-particle identity following from stationarity and Dyson's equation.\\
Implemented & Present in the inspected \code{dmft2} source.  It may still inherit the frozen-augmentation and finite-basis approximations of the underlying LAPW stress.\\
Proposed diagnostic & A controlled comparison that has not yet been added to the source.\\
Open audit item & A difference between the exact target and the present implementation that still requires numerical closure.\\
\bottomrule
\end{longtable}

The note uses a fixed band-index DMFT window.  The $\kb$-point weights are
constant under strain.  Muffin-tin radii, radial projector files, and the outer
radius of projector 5 are held fixed in bohr.  The base implementation uses
fixed local correlated axes.  A co-moving-axis correction is a later extension
for a self-energy constrained to remain diagonal in a local frame.

\section{Notation and dimensions}

\subsection{Indices and spaces}

\begin{longtable}{>{$}p{0.16\linewidth}<{$}p{0.76\linewidth}}
\toprule
\kb & Brillouin-zone point; $w_{\kb}$ is its fixed normalized integration weight.\\
n,m & Kohn--Sham band indices inside the fixed DMFT window $W_{\kb}$.\\
a,b & Bands stored by LAPW1 outside or inside $W_{\kb}$; $a\notin W_{\kb}$ denotes a complementary band.\\
\ell,\ell' & Composite correlated-orbital indices, including atom, shell, orbital, and spin block as needed.\\
\nu & One of the six independent symmetric strain channels, ordered as $(xx,yy,zz,yz,xz,xy)$.\\
\alpha,\beta & Cartesian indices in $\{x,y,z\}$; $\nu\equiv(\alpha\beta)$.\\
i\omega_j & Fermionic Matsubara frequency, $\omega_j=(2j+1)\pi/\beta$.\\
\beta & Inverse temperature; the code uses Rydberg energy units.\\
\Omega & Unstrained unit-cell volume.\\
\bottomrule
\end{longtable}

Let $N_b(\kb)$ be the number of bands in the DMFT window and $N_c$ the number
of active correlated orbitals.  The central matrices have dimensions

\begin{align}
 \rho_{\kb},\rho_{E,\kb},\epsilon_{\kb},G_{\kb} &\in
     \mathbb C^{N_b\times N_b},\\
 P_{\kb},U_{\kb} &\in \mathbb C^{N_b\times N_c},\\
 \Sig, S, X &\in \mathbb C^{N_c\times N_c},\\
 C_{\kb}&\in\mathbb C^{N_c\times N_b}.
\end{align}

The trace convention is explicit.  For example,

\begin{equation}
 \Tr_c(C_{\kb}U_{\kb,\nu})
 =\sum_{\ell n}(C_{\kb})_{\ell n}(U_{\kb,\nu})_{n\ell}.
 \label{eq:traceconvention}
\end{equation}

No dagger belongs in Eq.~\eqref{eq:traceconvention}; $C$ is the rectangular
one-particle gradient with the orientation chosen in Sec.~\ref{sec:C}.

\subsection{Strain and stress convention}

With deformation gradient

\begin{equation}
 F=I+\varepsilon,\qquad \mathbf r'=F\mathbf r,
 \qquad \mathbf G'=F^{-T}\mathbf G,
\end{equation}

the thermodynamic stress convention used in this note is

\begin{equation}
 \sigma_{\alpha\beta}=\frac{1}{\Omega}
 \left.\frac{\partial\Gamma}{\partial\varepsilon_{\alpha\beta}}\right|_0,
 \qquad P=-\frac13\Tr\sigma.
 \label{eq:stressdef}
\end{equation}

For an independent symmetric shear variable,

\begin{equation}
 \frac{\partial\varepsilon_{\gamma\delta}}
 {\partial\varepsilon_{\alpha\beta}}
 =\frac12(\delta_{\gamma\alpha}\delta_{\delta\beta}
          +\delta_{\gamma\beta}\delta_{\delta\alpha}).
\end{equation}

Consequently, real-space and reciprocal-space vectors obey

\begin{equation}
 \partial_{\alpha\beta}V_\gamma
 =\pm\frac12(\delta_{\gamma\alpha}V_\beta
             +\delta_{\gamma\beta}V_\alpha),
\end{equation}

with $+$ for real-space vectors and $-$ for reciprocal vectors.  A shear
finite difference must use the same convention or it will differ by a factor
of two.

\section{Bird's-eye stationary-functional derivation}

\subsection{Why the Klein functional is differentiated}

Write the stationary DFT+eDMFT functional schematically as

\begin{equation}
 \Gamma=\Gamma_{\rm DFT}[\rho]
 +\Gamma_{\rm band}[G;H,O,\Sigma^{\rm emb}]
 +\sum_R\left(\Phi_{{\rm imp},R}[G_R]-\Phi_{{\rm DC},R}[G_R]\right).
 \label{eq:schematicKlein}
\end{equation}

The detailed representation of individual double-counting and impurity terms
is not needed for the first geometric derivative.  At self-consistency,

\begin{equation}
 \frac{\delta\Gamma}{\delta G}=0,
 \qquad \frac{\delta\Gamma}{\delta\Sigma}=0,
 \qquad \frac{\delta\Gamma}{\delta\rho}=0.
\end{equation}

Therefore the total derivative with respect to a strain channel $\nu$ reduces
to the explicit derivative,

\begin{equation}
 \frac{\dd\Gamma}{\dd\varepsilon_\nu}
 =\left.\frac{\partial\Gamma}{\partial\varepsilon_\nu}
 \right|_{G,\Sigma,\rho}.
 \label{eq:stationarity}
\end{equation}

Terms containing $G_{,\nu}$, $\Sigma_{,\nu}$, and the impurity vertex cancel.
This is a consequence of stationarity, not a frozen-self-energy
approximation.  Explicit changes of $H$, $O$, and the embedding projector do
not cancel.

\subsection{Nonorthogonal LAPW form}

In an arbitrary nonorthogonal LAPW coefficient basis $\{\chi_K\}$,

\begin{equation}
 G_\chi^{-1}(i\omega)
 =(i\omega+\mu)O-H_0-\Sigma_\chi^{\rm emb}(i\omega).
\end{equation}

Define the frequency-summed coefficient-space matrices

\begin{align}
 \rho_\chi &=T\sum_jG_\chi(i\omega_j)e^{i\omega_j0^+},\\
 \rho_{E,\chi}&=T\sum_j(i\omega_j+\mu)
 G_\chi(i\omega_j)e^{i\omega_j0^+}.
\end{align}

The stationary one-particle derivative is

\begin{equation}
 D_\nu^{1p}
 =\sum_{\kb}\Ree\Tr\left[
 \rho_{\chi,\kb}H_{\kb,\nu}
 -\rho_{E,\chi,\kb}O_{\kb,\nu}\right]
 +D_\nu^{\rm emb}.
 \label{eq:nonorthogonal}
\end{equation}

The remaining DFT density, Hartree, exchange-correlation, core, nuclear, and
electrostatic terms assemble exactly as in the validated DFT stress.  The new
term $D_\nu^{\rm emb}$ is the explicit derivative of the correlated embedding
at fixed $\Sig$.

\identitybox{Equation~\eqref{eq:nonorthogonal} is the organizing equation.
It contains a paired $H_{,\nu}$ and $O_{,\nu}$ contraction.  Their individual
magnitudes can be very large, so numerical or analytic replacements must
preserve the stationary pair.}

\section{Green functions, density matrices, and energy moments}

In the orthonormal Kohn--Sham band basis of the DMFT window,

\begin{align}
 \Sigma_{\kb}^{\rm emb}(i\omega)
 &=U_{\kb}\Sig(i\omega)U_{\kb}^\dagger,\\
 G_{\kb}^{-1}(i\omega)
 &=(i\omega+\mu)I-\epsilon_{\kb}
 -\Sigma_{\kb}^{\rm emb}(i\omega),
 \label{eq:dyson}
\end{align}

where $\Sig=\Sigma_{\rm imp}-V_{\rm DC}$ and $\epsilon_{\kb}$ is diagonal.
The code constructs Hermitian frequency-summed matrices

\begin{align}
 \rho_{\kb}&=T\sum_jG_{\kb}(i\omega_j)e^{i\omega_j0^+},\\
 \rho_{E,\kb}&=T\sum_j(i\omega_j+\mu)
 G_{\kb}(i\omega_j)e^{i\omega_j0^+}.
 \label{eq:rhorhoE}
\end{align}

Dyson's equation gives the left and right identities

\begin{align}
 B_{L,\kb}&\equiv\rho_{E,\kb}-\epsilon_{\kb}\rho_{\kb}
 =T\sum_j\Sigma_{\kb}^{\rm emb}(i\omega_j)G_{\kb}(i\omega_j),\\
 B_{R,\kb}&\equiv\rho_{E,\kb}-\rho_{\kb}\epsilon_{\kb}
 =B_{L,\kb}^\dagger.
 \label{eq:BLBR}
\end{align}

The Hermitian self-energy moment is

\begin{equation}
 M_{\Sigma,\kb}
 \equiv\rho_{E,\kb}-\frac12\{\epsilon_{\kb},\rho_{\kb}\}
 =\frac12(B_{L,\kb}+B_{R,\kb}).
 \label{eq:Msigma}
\end{equation}

Using Eq.~\eqref{eq:Msigma}, the one-particle LAPW derivative can be separated
into a DFT-like part and the genuinely new energy-moment metric part,

\begin{align}
 D_\nu^{1p}
 &=D_\nu^{\rm DFT-like}
 -\sum_{\kb}\Ree\Tr[M_{\Sigma,\kb}O_{\kb,\nu}],\\
 D_\nu^{\rm DFT-like}
 &=\sum_{\kb}\Ree\Tr\left[
 \rho_{\kb}H_{\kb,\nu}
 -\frac12\{\epsilon_{\kb},\rho_{\kb}\}O_{\kb,\nu}
 \right].
 \label{eq:dftlike}
\end{align}

\warningbox{The $M_\Sigma O_{,\nu}$ term requires the complete generalized
APW+LO metric derivative.  The current implementation obtains an overlap
kernel from a selected rearranged valence-correction kernel.  Direct matrix
tests show that this isolated analytic channel does not yet reproduce the full
metric derivative in a material-independent way.}

\section{Generalized LAPW strain matrices}
\label{sec:TBO}

\subsection{Fixed-coefficient basis response}

Let

\begin{equation}
 |\psi_{n\kb}\rangle=\sum_a c_{an\kb}|\phi_{a\kb}\rangle
\end{equation}

be a Kohn--Sham state expanded in the LAPW+LO basis.  Its fixed-coefficient
basis response is

\begin{equation}
 |\psi^{B,\nu}_{n\kb}\rangle
 =\sum_a c_{an\kb}|\partial_\nu\phi_{a\kb}\rangle.
 \label{eq:basisresponse}
\end{equation}

The superscript $B$ means ``basis response'': the eigenvector coefficients
$c_{an\kb}$ are held fixed.  Eigenvector rotation into complementary bands is
treated separately in Sec.~\ref{sec:window}.

Define one-sided matrices

\begin{align}
 b^\nu_{mn}(\kb)&=
 \langle\psi^{B,\nu}_{m\kb}|H|\psi_{n\kb}\rangle,\\
 o^\nu_{mn}(\kb)&=
 \langle\psi^{B,\nu}_{m\kb}|\psi_{n\kb}\rangle.
 \label{eq:one-sided}
\end{align}

Their physical Hermitian polarizations are

\begin{align}
 \mathcal B^\nu&=b^\nu+b^{\nu\dagger},\\
 \mathcal O^\nu&=o^\nu+o^{\nu\dagger}.
 \label{eq:BOherm}
\end{align}

Thus

\begin{align}
 \mathcal B^\nu_{mn}&=
 \langle\psi^{B,\nu}_{m}|H|\psi_n\rangle
 +\langle\psi_m|H|\psi^{B,\nu}_n\rangle,\\
 \mathcal O^\nu_{mn}&=
 \langle\psi^{B,\nu}_{m}|\psi_n\rangle
 +\langle\psi_m|\psi^{B,\nu}_n\rangle.
\end{align}

For $m=n$,

\begin{equation}
 \mathcal B^\nu_{nn}-\epsilon_n\mathcal O^\nu_{nn}
 =2\Ree\langle\psi^{B,\nu}_n|(H-\epsilon_n)|\psi_n\rangle,
 \label{eq:Blaha20general}
\end{equation}

which is the diagonal expression in Blaha et al., Eq.~(20).  Blaha does not
write the full $m\ne n$ matrices because ordinary DFT occupations are diagonal
in the Kohn--Sham basis.

\subsection{Kinetic strain matrix}

For $\nu=(\alpha\beta)$, the generalized valence kinetic matrix is

\begin{equation}
 T^{\alpha\beta}_{mn}(\kb)
 =\frac12\int_\Omega\dd^3r\,
 \psi_{m\kb}^*(\mathbf r)
 (\partial_\alpha\partial_\beta+
  \partial_\beta\partial_\alpha)
 \psi_{n\kb}(\mathbf r).
 \label{eq:Tmn}
\end{equation}

Its diagonal contraction is Blaha et al., Eq.~(19).  Appendix G, Eq.~(G1),
provides the expanded LAPW expression.  A complete implementation includes
the muffin-tin term, the interstitial term, the APW boundary/surface term, and
the selected scalar-relativistic replacement.

\subsection{The energy-linear correction kernel}

The WIEN2k routine \code{VAL\_CORR\_STRESS} evaluates a one-sided quadratic
kernel that is linear in a supplied scalar energy $e$,

\begin{equation}
 K^\nu(e)=h^\nu-e\,q^\nu.
 \label{eq:Klinear}
\end{equation}

Here $h^\nu$ and $q^\nu$ denote the routine's affine-strain kernels.  Their
sign relative to Eqs.~\eqref{eq:one-sided}--\eqref{eq:BOherm} is fixed by the
WIEN2k generator convention and by the DFT reduction.  The source accumulates

\begin{equation}
 D_\nu^{\rm corr,impl}
 =-2\Ree\Tr(\rho h^\nu)
 +2\Ree\Tr(\rho_E q^\nu).
 \label{eq:implcorr}
\end{equation}

The two required pieces are extracted without rewriting hundreds of angular
terms:

\begin{equation}
 K^\nu(0)=h^\nu,
 \qquad K^\nu(0)-K^\nu(1)=q^\nu.
 \label{eq:Ksplit}
\end{equation}

Blaha's fully expanded diagonal correction is Appendix G, Eq.~(G2).  The
off-diagonal generalization is not separately printed in that paper; it is
obtained by Hermitian polarization as above.

\section{Why pseudo-band diagonalization gives the off-diagonal result}

Let $A$ be any Hermitian band-space matrix and diagonalize a Hermitian density
matrix,

\begin{equation}
 \rho=VfV^\dagger,\qquad f=\diag(f_p).
\end{equation}

Define pseudo-states

\begin{equation}
 |\chi_p\rangle=\sum_nV_{np}|\psi_n\rangle.
\end{equation}

Then

\begin{equation}
 \Tr(\rho A)=\sum_pf_p\langle\chi_p|A|\chi_p\rangle.
 \label{eq:pseudobandidentity}
\end{equation}

Equation~\eqref{eq:pseudobandidentity} is exact and includes every
off-diagonal $m\ne n$ term.  It does not assume that $\rho$ is positive;
signed eigenvalues are allowed.  The code applies this identity twice:

\begin{align}
 \rho&=VfV^\dagger &&\text{for }T^\nu\text{ and }h^\nu,\\
 \rho_E&=WgW^\dagger &&\text{for }q^\nu.
\end{align}

For a pseudo-state with vector $v_p$,

\begin{align}
 k_p^\nu(0)&=v_p^\dagger h^\nu v_p,\\
 k_p^\nu(0)-k_p^\nu(1)&=v_p^\dagger q^\nu v_p.
\end{align}

Therefore the diagonal WIEN2k quadratic kernel acts as a matrix-free
implementation of the generalized matrices.

\implemented{The current code does not store dense arrays $T^\nu_{mn}$,
$\mathcal B^\nu_{mn}$, or $\mathcal O^\nu_{mn}$.  It implements their Hermitian
contractions by pseudo-band polarization in \code{cmp\_dmft\_stress.f90}.}

\section{Exact mapping to the present stress implementation}

The present matrix-free path is:

\begin{enumerate}
\item Diagonalize $\rho$ and rotate LAPW $A/B/C$ coefficients to its
      pseudo-band eigenvectors.
\item Evaluate muffin-tin kinetic stress with \code{val\_kinstrmt}.
\item Add the scalar-relativistic replacement with
      \code{val\_rela\_kinstrmt} when \code{CALC=RELA}.
\item Evaluate the APW kinetic surface term with
      \code{stress\_APW\_surface}.
\item Set the energy argument to zero and call
      \code{val\_corr\_stress}; this gives $h^\nu$.
\item Evaluate interstitial kinetic stress once per spin and $\kb$ point.
\item Diagonalize $\rho_E$, rotate coefficients to its pseudo-bands, and call
      \code{val\_corr\_stress} at $e=0$ and $e=1$; the difference gives
      $q^\nu$.
\item Rotate equivalent-atom tensors to global Cartesian axes, reduce over
      MPI ranks, and apply crystal symmetrization.
\end{enumerate}

The non-spherical correction is decomposed into matching, radial, and angular
parts.  The current conservative default uses the trace combination; a
literal full affine form is available as a diagnostic.  This choice is part
of the DFT stress audit and is not a new DMFT term.

\warningbox{``Implemented off-diagonal matrices'' means that the contractions
needed by Hermitian $\rho$ and $\rho_E$ are implemented.  The code cannot
currently print an arbitrary individual matrix element $T^\nu_{mn}$,
$B^\nu_{mn}$, or $O^\nu_{mn}$ without a new matrix-assembly diagnostic.}

\section{Raw and normalized correlated projectors}

\subsection{Definitions}

Let the raw finite-window projector be

\begin{equation}
 (P_{\kb})_{n\ell}=\langle\psi_{n\kb}|\varphi_\ell\rangle,
 \qquad P_{\kb}\in\mathbb C^{N_b\times N_c}.
 \label{eq:rawP}
\end{equation}

The global correlated-orbital overlap and its inverse square root are

\begin{equation}
 S=\sum_{\kb}w_{\kb}P_{\kb}^\dagger P_{\kb},
 \qquad X=S^{-1/2}.
 \label{eq:SX}
\end{equation}

The projector used by the lattice problem is

\begin{equation}
 \boxed{U_{\kb}=P_{\kb}X.}
 \label{eq:normalizedU}
\end{equation}

It obeys the global normalization

\begin{equation}
 \sum_{\kb}w_{\kb}U_{\kb}^\dagger U_{\kb}=I_c,
 \label{eq:globalnorm}
\end{equation}

but generally not $U_{\kb}^\dagger U_{\kb}=I_c$ at each $\kb$.

\subsection{Complete derivative}

The fixed weights satisfy $w_{\kb,\nu}=0$.  Hence

\begin{align}
 U_{\kb,\nu}&=P_{\kb,\nu}X+P_{\kb}X_{,\nu},
 \label{eq:Uderiv}\\
 S_{,\nu}&=\sum_{\kb}w_{\kb}
 \left(P_{\kb,\nu}^\dagger P_{\kb}
 +P_{\kb}^\dagger P_{\kb,\nu}\right).
 \label{eq:Sderiv}
\end{align}

The derivative of an inverse square root is a matrix Fr\'echet derivative.
If $S=V\diag(s_i)V^\dagger$, then

\begin{equation}
 (V^\dagger X_{,\nu}V)_{ij}
 =f^{[1]}(s_i,s_j)(V^\dagger S_{,\nu}V)_{ij},
 \qquad f(s)=s^{-1/2},
 \label{eq:Frechet}
\end{equation}

where

\begin{equation}
 f^{[1]}(s_i,s_j)=
 \begin{cases}
 \dfrac{s_i^{-1/2}-s_j^{-1/2}}{s_i-s_j},&s_i\ne s_j,\\[6pt]
 -\dfrac12s_i^{-3/2},&s_i=s_j.
 \end{cases}
\end{equation}

Equivalently, with $Y=S^{1/2}=X^{-1}$,

\begin{equation}
 YX_{,\nu}+X_{,\nu}Y=-XS_{,\nu}X.
 \label{eq:sylvester}
\end{equation}

Differentiating Eq.~\eqref{eq:globalnorm} supplies a stringent check,

\begin{equation}
 \sum_{\kb}w_{\kb}
 \left(U_{\kb,\nu}^\dagger U_{\kb}
 +U_{\kb}^\dagger U_{\kb,\nu}\right)=0.
 \label{eq:normcheck}
\end{equation}

\implemented{The normalization branch $P_{\kb}X_{,\nu}$ is present through
the relative derivative $X^{-1}X_{,\nu}$ in
\code{AddProjectorStressNormalization}.  It is the stress analogue of the
normalization correction added to the DFT+eDMFT force.}

\subsection{What belongs to $P_{\kb,\nu}$}

It is useful to write

\begin{equation}
 P_{\kb,\nu}=P_{\kb,\nu}^{\rm geom}
 +P_{\kb,\nu}^{\rm radial}
 +P_{\kb,\nu}^{\rm subspace}
 +P_{\kb,\nu}^{\rm axes}.
 \label{eq:Psplit}
\end{equation}

The terms mean:

\begin{description}
\item[Geometry.] Reciprocal vectors, phase factors, angular functions,
matching coefficients, volume normalization, APW and LO contributions.
\item[Radial.] A change of the stored radial projector itself.  This is zero
in the base convention because the same \code{projectorw.dat}, $R_{\rm MT}$,
and outer radius are used at every strain.
\item[Subspace.] Rotation of the fixed band-index window into complementary
stored bands; Sec.~\ref{sec:window}.
\item[Axes.] Rotation of a geometry-following correlated frame.  It is zero
for the fixed-axis base implementation.
\end{description}

For a current spherical muffin-tin domain with fixed $R_{\rm MT}$, the
atom-centered coordinate $\mathbf x=\mathbf r-\boldsymbol\tau_a$ spans the same
fixed ball $|\mathbf x|<R_{\rm MT}$ at every strain.  Therefore no extra
Reynolds boundary term is generated merely because a radial projector is
finite at $R_{\rm MT}$.  For projector 5, the tail outside $R_{\rm MT}$ must be
included consistently; then $R_{\rm MT}$ is not a support boundary at all.

\section{The one-particle projector gradient $C_{\kb}$}
\label{sec:C}

At fixed $\Sig$, the embedded self-energy varies as

\begin{equation}
 \delta\Sigma_{\kb}^{\rm emb}
 =(\delta U_{\kb})\Sig U_{\kb}^\dagger
 +U_{\kb}\Sig(\delta U_{\kb})^\dagger.
\end{equation}

Define

\begin{equation}
 \boxed{
 C_{\kb}=\frac{w_{\kb}}{\beta}
 \sum_j\Sig(i\omega_j)U_{\kb}^\dagger G_{\kb}(i\omega_j)
 }
 \quad\in\mathbb C^{N_c\times N_b}.
 \label{eq:Cdirect}
\end{equation}

The explicit normalized-projector derivative is

\begin{equation}
 \boxed{
 D_\nu^{\rm proj}
 =2\Ree\sum_{\kb}\Tr_c(C_{\kb}U_{\kb,\nu}).
 }
 \label{eq:projectorstress}
\end{equation}

The factor $2\Ree$ accounts for the two Hermitian-conjugate branches.  If both
branches are explicitly constructed, it must not be applied a second time.

Using Eq.~\eqref{eq:BLBR},

\begin{equation}
 \boxed{U_{\kb}C_{\kb}=B_{L,\kb}
 =\rho_{E,\kb}-\epsilon_{\kb}\rho_{\kb}.}
 \label{eq:UC}
\end{equation}

This is a band-space identity.  $U_{\kb}C_{\kb}$ is
$N_b\times N_b$, while $C_{\kb}U_{\kb,\nu}$ is $N_c\times N_c$.

\section{SVD, pseudoinverse reconstruction, and what can be lost}

Let

\begin{equation}
 U_{\kb}=L_{\kb}s_{\kb}R_{\kb}^\dagger
\end{equation}

be a singular-value decomposition.  A small orbital-space dimension makes
$U^\dagger U$ inexpensive to invert, but does not guarantee it is nonsingular.
Symmetry or an insufficient band window can make one or more singular values
zero at an individual $\kb$.

If $U_{\kb}$ has full column rank,

\begin{equation}
 C_{\kb}=(U_{\kb}^\dagger U_{\kb})^{-1}
 U_{\kb}^\dagger B_{L,\kb}.
\end{equation}

The implemented general reconstruction uses the Moore--Penrose pseudoinverse,

\begin{equation}
 \boxed{
 C_{\kb}^{\rm MP}=U_{\kb}^{\rm MP}B_{L,\kb}
 =(U_{\kb}^\dagger U_{\kb})^{\rm MP}
 U_{\kb}^\dagger B_{L,\kb}.}
 \label{eq:CMP}
\end{equation}

The notation $A^{\rm MP}$ is used to avoid confusing a pseudoinverse with a
Hermitian conjugate $A^\dagger$.

If Eq.~\eqref{eq:UC} is exact for a directly accumulated $C^{\rm dir}$, then

\begin{equation}
 C_{\kb}^{\rm MP}
 =U_{\kb}^{\rm MP}U_{\kb}C_{\kb}^{\rm dir}
 \equiv\Pi_{\kb}C_{\kb}^{\rm dir},
\end{equation}

where $\Pi_{\kb}$ projects onto the local subspace represented by the chosen
band window.  Hence

\begin{align}
 C_{\kb}^{\rm dir}&=C_{\kb}^{\rm MP}+C_{\kb}^{\perp},\\
 U_{\kb}C_{\kb}^{\perp}&=0.
\end{align}

Although $C^\perp$ is invisible in $B_L=UC$, it can contribute to the stress:

\begin{equation}
 \Delta D_\nu^{\rm proj}
 =2\Ree\Tr(C_{\kb}^{\perp}U_{\kb,\nu}).
 \label{eq:nullstress}
\end{equation}

This is the precise reason the equation-of-motion reconstruction can be
insufficient at a rank-deficient $\kb$ point.

\implemented{The current \code{BuildWindowResponseMatrices} diagonalizes the
small Gram matrix $U_W^\dagger U_W$, applies a relative eigenvalue cutoff, and
constructs Eq.~\eqref{eq:CMP}.  It reports the reconstruction residual when
projector diagnostics are enabled.}

\section{Fixed-index window and complementary-band response}
\label{sec:window}

The DMFT window membership is fixed by band index, so there is no derivative
of a sharp energy cutoff.  The selected eigenvectors nevertheless rotate into
bands outside the window.  For a nondegenerate window state $n$,

\begin{equation}
 |\psi_{n,\nu}\rangle_\perp
 =\sum_{a\notin W}|\psi_a\rangle
 \frac{\langle\psi_a|H_{,\nu}-\epsilon_nO_{,\nu}|\psi_n\rangle}
 {\epsilon_n-\epsilon_a}.
 \label{eq:sternheimerbands}
\end{equation}

The exact sum contains all complementary bands.  The implementation uses the
additional bands already present in \code{case.vector} or
\code{case.vectorso}; convergence can be improved by asking LAPW1 to write
more bands.

Let $U_{\rm full}$ contain projector rows for every stored band and let

\begin{equation}
 K_{\rm full}=U_{\rm full}C^{\rm MP}.
\end{equation}

For $a\notin W$ and $n\in W$, the code builds auxiliary Hermitian response
matrices

\begin{align}
 (\rho^{\rm resp})_{an}
 &=\frac{(K_{\rm full})_{an}}{\epsilon_n-\epsilon_a},
 & (\rho^{\rm resp})_{na}&=(\rho^{\rm resp})_{an}^*,\\
 (\rho_E^{\rm resp})_{an}
 &=\epsilon_n(\rho^{\rm resp})_{an},
 & (\rho_E^{\rm resp})_{na}&=(\rho_E^{\rm resp})_{an}^*.
 \label{eq:responsematrices}
\end{align}

The corresponding response is evaluated by reusing the LAPW kernels,

\begin{equation}
 D_\nu^{\rm resp}
 =\Tr(\rho^{\rm resp}H_{,\nu})
 -\Tr(\rho_E^{\rm resp}O_{,\nu}),
 \label{eq:windowresponse}
\end{equation}

with the source's calibrated sign convention.  The response matrices are a
contraction device; they must not be added to the physical charge density.
Near cross-window degeneracies, ordinary denominators are not reliable and a
degenerate spectral-projector treatment is ultimately preferable.

\section{Direct $C_{\kb}$ accumulation: proposed comparison switch}

The pseudoinverse can be eliminated by accumulating Eq.~\eqref{eq:Cdirect}
inside the existing Matsubara loop.  At a positive Matsubara frequency the
paired contribution is

\begin{equation}
 \Sig(i\omega)U^\dagger G(i\omega)
 +\Sig^\dagger(i\omega)U^\dagger G^\dagger(i\omega).
 \label{eq:Cpair}
\end{equation}

Because $C$ is rectangular and generally non-Hermitian, it must not be
Hermitian-symmetrized.  The negative-frequency partner in
Eq.~\eqref{eq:Cpair} is not $C(i\omega)^\dagger$; that object has transposed
dimensions and a different matrix order.

\proposed{Add
\code{EDMFT\_STRESS\_C\_MODE=PINV|DIRECT|COMPARE}.  \code{PINV} retains the
current production path.  \code{DIRECT} uses Eq.~\eqref{eq:Cdirect}.
\code{COMPARE} computes both while initially leaving the production result
unchanged.}

The comparison should report

\begin{align}
 r_C&=\frac{\norm{C^{\rm dir}-C^{\rm MP}}_F}
 {\max(\norm{C^{\rm dir}}_F,\epsilon_{\rm num})},\\
 r_B^{\rm dir}&=\frac{\norm{UC^{\rm dir}-B_L}_F}
 {\max(\norm{B_L}_F,\epsilon_{\rm num})},\\
 r_B^{\rm MP}&=\frac{\norm{UC^{\rm MP}-B_L}_F}
 {\max(\norm{B_L}_F,\epsilon_{\rm num})},\\
 \Delta\sigma_\nu^{\rm proj}
 &=\frac{2}{\Omega}\Ree\Tr\left[
 (C^{\rm dir}-C^{\rm MP})U_{,\nu}\right].
 \label{eq:Cdiagnostics}
\end{align}

Three band-space objects should also be compared independently:

\begin{enumerate}
\item $B_L^{\rm EOM}=\rho_E-\epsilon\rho$;
\item $B_L^{\rm band}=T\sum\Sigma^{\rm emb}G$, already accumulated by a
      diagnostic in \code{cmp\_dmft\_weights2};
\item $B_L^{\rm local}=UC^{\rm dir}$.
\end{enumerate}

This three-way test separates an energy-moment/tail problem from a local
projector mapping problem and from information lost by the pseudoinverse.

\section{Matsubara summation used by $\rho_E$ and the proposed direct $C$}

The implementation uses two frequency regions:

\begin{enumerate}
\item the first consecutive Matsubara frequencies are summed explicitly;
\item the high-frequency region is stored on a sparse logarithmic mesh and
      summed with precomputed cardinal cubic-spline weights
      \code{matsum\_weight}.
\end{enumerate}

The nonlinear resolvent is formed first; the final matrix-valued integrand is
then multiplied by the cubic-spline quadrature weight.  This is more accurate
than interpolating poles or eigenvalues before forming the resolvent.

For the dynamic self-energy tail,

\begin{equation}
 \Sig(i\omega)-\Sigma_\infty
 =\frac{\Sigma_1}{i\omega}
 +\frac{\Sigma_2}{(i\omega)^2}+\cdots.
\end{equation}

The code adaptively estimates $\Sigma_1$: a smooth tail uses a two-moment fit,
a noisy tail uses a short plateau average, and a clearly bad or increasing
tail falls back to the robust endpoint estimate.  In the energy-moment sum,
the leading $\Sigma_1/(i\omega)^2$ contribution is subtracted from the sparse
cubic quadrature and restored over exactly the same frequency interval.

If the positive-frequency sparse tail starts at fermionic index $j_0$ with
$\omega_j=(2j-1)\pi/\beta$, the scalar used by the code is

\begin{equation}
 M_2^+(j_0)=\sum_{j=j_0}^{\infty}\frac{1}{\omega_j^2}
 =\frac{\beta^2}{8}
 -\sum_{j=1}^{j_0-1}
 \left[\frac{\beta}{(2j-1)\pi}\right]^2.
 \label{eq:tailMoment}
\end{equation}

The final paired Matsubara factor supplies $2/\beta$.  Equation
\eqref{eq:tailMoment} is safe because it is applied only to the high-frequency
tail after the same term has been removed from the numerical residual; it is
not added as a singular approximation at low frequency.

The direct-$C$ accumulator must use exactly the same cubic weights, positive
and negative matrix ordering, $\Sigma_1$ subtraction, and analytic restoration
as $B_L$.  Otherwise $UC^{\rm dir}=B_L$ would fail for a purely numerical
reason.

\section{Standard DFT stress channels retained in eDMFT}

The electronic DMFT generalization does not replace the remaining DFT stress
machinery.  The full result retains:

\begin{itemize}
\item valence kinetic stress, including muffin-tin, interstitial, APW surface,
      and relativistic corrections;
\item valence incomplete-basis/Pulay correction;
\item core correction from \code{lcore};
\item electrostatic and nuclear contributions from \code{lapw0};
\item exchange-correlation stress, including the GGA contribution where
      applicable;
\item spin-channel assembly and final crystal symmetrization;
\item the explicit normalized-projector term and fixed-window response.
\end{itemize}

The eDMFT matrices $\rho$ and $\rho_E$ replace diagonal DFT occupations only
inside the valence one-particle contractions.  A naive substitution of
diagonal occupations loses off-diagonal coherence and gives the overlap term
the wrong energy weight.

\section{Source-level equation map}

\begin{longtable}{p{0.28\linewidth}p{0.31\linewidth}p{0.33\linewidth}}
\toprule
Mathematical object & Current source location & Implementation strategy\\
\midrule
$\rho,\rho_E$ & \code{cmp\_dmft\_weights.f90} & Matsubara sum, cubic sparse-tail quadrature, analytic $\Sigma_1$ restoration.\\
$T^\nu$ contraction & \code{cmp\_dmft\_stress.f90} & $\rho$ pseudo-bands passed to muffin-tin, interstitial, APW surface, and RELA kernels.\\
$h^\nu$ contraction & \code{cmp\_dmft\_stress.f90} & \code{val\_corr\_stress} evaluated at $e=0$.\\
$q^\nu$ contraction & \code{cmp\_dmft\_stress.f90} & $\rho_E$ pseudo-bands; $K(0)-K(1)$.\\
$U_{,\nu}$ geometry & \code{l2main.f90}, projector strain routines & Fixed-coefficient raw-projector derivative.\\
$PX_{,\nu}$ & \texttt{AddProjectorStress-}\newline\texttt{Normalization} & Relative inverse-square-root derivative.\\
$C^{\rm MP}$ & \texttt{BuildWindowResponse-}\newline\texttt{Matrices} & Eigenvalue-cutoff Moore--Penrose inverse of $U_W^\dagger U_W$.\\
Complementary bands & \texttt{BuildWindowResponse-}\newline\texttt{Matrices} & Eqs.~\eqref{eq:responsematrices}--\eqref{eq:windowresponse} using stored LAPW1 bands.\\
$T\sum\Sigma^{\rm emb}G$ check & \code{cmp\_dmft\_weights.f90} & Direct band-space diagnostic with the same Matsubara tail.\\
$C^{\rm dir}$ & not yet present & Proposed local-by-band accumulation, Eq.~\eqref{eq:Cdirect}.\\
\bottomrule
\end{longtable}

\section{What has been established and what remains open}

\subsection{Established algebraic facts}

\begin{itemize}
\item The first derivative requires no impurity two-particle vertex.
\item $\rho_E$, not $\epsilon\rho$, multiplies the generalized-overlap
      derivative.
\item Blaha's diagonal Eqs.~(19) and (20) generalize uniquely to the Hermitian
      bilinear matrices in Sec.~\ref{sec:TBO}.
\item Pseudo-band diagonalization evaluates the required off-diagonal
      contractions exactly for Hermitian $\rho$ and $\rho_E$.
\item The projector must be differentiated after normalization:
      $U_{,\nu}=P_{,\nu}X+PX_{,\nu}$.
\item A fixed band-index window still has a complementary-band eigenvector
      response.
\item $C^{\rm MP}$ can discard a null-space component that contributes through
      $U_{,\nu}$.
\end{itemize}

\subsection{Current numerical audit state}

Direct finite differences of complete LAPW $H$ and $O$ matrices satisfy the
generalized eigenvalue checks and reproduce finite differences of DFT band
energies.  This validates the operator-level paired expression
$\Tr(\rho H_{,\nu})-\Tr(\rho_EO_{,\nu})$.  The principal unresolved problems
are:

\begin{enumerate}
\item the current analytic isolated $M_\Sigma O_{,\nu}$ channel does not equal
      the complete metric derivative by one universal sign or factor;
\item DFT baseline residuals are not yet uniformly below the desired
      sub-
      \SI{0.1}{GPa} scale in every material;
\item the effect of $C^\perp$ at singular $\kb$ points has not been measured
      because $C^{\rm dir}$ is not yet accumulated;
\item RELA operator-level closure, shear components, and magnetic two-spin
      assembly remain later validation stages.
\end{enumerate}

The recommended nonmagnetic comparison sequence is SrVO$_3$, MnO, FeSe NREL,
and FeSe RELA; magnetic MnTe should follow after the nonmagnetic identities
close.

\section{Validation protocol}

\subsection{Algebraic checks at every $\kb$}

\begin{align}
 \delta_O&=\norm{A^\dagger OA-I}_F,\\
 \delta_H&=\norm{A^\dagger HA-\epsilon}_F,\\
 \delta_E&=\norm{\rho_E-\tfrac12\{\epsilon,\rho\}-M_\Sigma}_F,\\
 \delta_N&=\norm{\sum_\kb w_\kb(U_{,\nu}^\dagger U+U^\dagger U_{,\nu})}_F,\\
 \delta_C&=\norm{UC-B_L}_F.
\end{align}

The singular spectrum of $U_\kb^\dagger U_\kb$ should be printed together
with all pseudoinverse cutoffs.

\subsection{Finite-difference checks}

For a central step $h$,

\begin{equation}
 D_\nu^{\rm FD}(h)
 =\frac{\Gamma(+h)-\Gamma(-h)}{2h}.
\end{equation}

Use at least two steps and, where smooth, Richardson extrapolation,

\begin{equation}
 D_\nu^{\rm Rich}
 =\frac{4D_\nu^{\rm FD}(h/2)-D_\nu^{\rm FD}(h)}{3}.
\end{equation}

The $-h$, $0$, and $+h$ calculations must have identical $R_{\rm MT}$,
projector radial data, band-index window, number of stored bands, $\kb$ mesh,
and reciprocal cutoffs.  At frozen $\Sigma$, the DFT charge must be relaxed.
For a fully stationary Klein comparison the impurity functional difference
$\Phi_{\rm imp}-\Phi_{\rm DC}$ must be included consistently; its omission is
not an analytic basis error.

\section{Compact formula sheet}

\begin{align}
 \sigma_\nu&=\Omega^{-1}\partial_\nu\Gamma,\\
 G_\kb^{-1}&=(i\omega+\mu)I-\epsilon_\kb-U_\kb\Sig U_\kb^\dagger,\\
 \rho_\kb&=T\sum_\omega G_\kb,
 &\rho_{E,\kb}&=T\sum_\omega(i\omega+\mu)G_\kb,\\
 M_{\Sigma,\kb}&=\rho_{E,\kb}-\tfrac12\{\epsilon_\kb,\rho_\kb\},\\
 P_\kb&:N_b\times N_c,
 &S&=\sum_\kb w_\kb P_\kb^\dagger P_\kb,\\
 X&=S^{-1/2},
 &U_\kb&=P_\kb X,\\
 U_{\kb,\nu}&=P_{\kb,\nu}X+P_\kb X_{,\nu},\\
 C_\kb&=\frac{w_\kb}{\beta}\sum_\omega
 \Sig U_\kb^\dagger G_\kb,
 &D_\nu^{\rm proj}&=2\Ree\sum_\kb\Tr(C_\kb U_{\kb,\nu}),\\
 B_{L,\kb}&=\rho_{E,\kb}-\epsilon_\kb\rho_\kb=U_\kb C_\kb,\\
 C_\kb^{\rm MP}&=(U_\kb^\dagger U_\kb)^{\rm MP}
 U_\kb^\dagger B_{L,\kb},\\
 T_{mn}^{\alpha\beta}&=\tfrac12\langle\psi_m|
 \partial_\alpha\partial_\beta+partial_\beta\partial_\alpha|\psi_n\rangle,\\
 \mathcal B^\nu&=b^\nu+b^{\nu\dagger},
 &\mathcal O^\nu&=o^\nu+o^{\nu\dagger},\\
 b_{mn}^\nu&=\langle\psi_m^{B,\nu}|H|\psi_n\rangle,
 &o_{mn}^\nu&=\langle\psi_m^{B,\nu}|\psi_n\rangle,\\
 K^\nu(e)&=h^\nu-eq^\nu,
 &q^\nu&=K^\nu(0)-K^\nu(1).
\end{align}

\section{Recommended next implementation step}

The next source change should be narrowly diagnostic:

\begin{enumerate}
\item accumulate $C_\kb^{\rm dir}$ inside \code{cmp\_dmft\_weights2}, reusing
      the existing $G$ and $\Sigma$ at no additional matrix inversion cost;
\item apply the identical cubic Matsubara weights and high-frequency tail
      restoration used for $B_L$;
\item add \code{PINV}, \code{DIRECT}, and \code{COMPARE} modes;
\item in \code{COMPARE}, print Eqs.~\eqref{eq:Cdiagnostics} and both projector
      stress tensors without altering the default production answer;
\item run the same executable on SrVO$_3$, MnO, FeSe NREL/RELA, then MnTe.
\end{enumerate}

If $C^{\rm dir}$ and $C^{\rm MP}$ agree in stress despite singular values,
the remaining error lies in the LAPW $H/O$ kernels or another stress channel.
If they differ by the observed numerical residual, Eq.~\eqref{eq:nullstress}
identifies the missing projector-gradient component directly.

\clearpage
\appendix

\section{Derivation of the inverse-square-root Sylvester equation}

The normalized metric satisfies

\begin{equation}
 XSX=I.
\end{equation}

Differentiating gives

\begin{equation}
 X_{,\nu}SX+XS_{,\nu}X+XSX_{,\nu}=0.
\end{equation}

Since $SX=XS=S^{1/2}=Y$,

\begin{equation}
 X_{,\nu}Y+YX_{,\nu}=-XS_{,\nu}X,
\end{equation}

which is Eq.~\eqref{eq:sylvester}.  The solution is unique when $S$ is
positive definite.  If the global $S$ itself is singular, the requested
correlated set is not represented by the total selected band manifold and an
SVD/truncated-subspace definition is required before differentiating.

\section{Polarization proof for the generalized kernels}

For $\rho=VfV^\dagger$,

\begin{align}
 \Tr(\rho A)
 &=\Tr(VfV^\dagger A)
 =\Tr(fV^\dagger AV)\\
 &=\sum_pf_p(V^\dagger AV)_{pp}
 =\sum_pf_pv_p^\dagger Av_p.
\end{align}

Thus a routine capable of evaluating $v^\dagger Av$ for an arbitrary linear
combination of Kohn--Sham states evaluates the full off-diagonal contraction.
The same argument applies independently to $\rho_E$, even when its eigenvalues
are negative.

\section{Relation to Blaha's numbered equations}

\begin{longtable}{p{0.18\linewidth}p{0.34\linewidth}p{0.39\linewidth}}
\toprule
Blaha equation & Published object & Present generalization\\
\midrule
Eq.~(19) & Diagonal valence kinetic stress $T_{nn}^{\alpha\beta}$ & Hermitian matrix $T_{mn}^{\alpha\beta}$, Eq.~\eqref{eq:Tmn}, contracted with off-diagonal $\rho$.\\
Eq.~(20) & Diagonal IBS/Pulay form $2\Ree\langle\psi_n^{B,\nu}|H-\epsilon_n|\psi_n\rangle$ & Matrices $\mathcal B^\nu$ and $\mathcal O^\nu$, Eqs.~\eqref{eq:one-sided}--\eqref{eq:Blaha20general}.\\
Eq.~(G1) & Expanded diagonal kinetic LAPW formula & Reused on $\rho$ pseudo-bands.\\
Eq.~(G2) & Expanded diagonal valence-correction formula & Reused on $\rho$ and $\rho_E$ pseudo-bands through $K(0)$ and $K(0)-K(1)$.\\
\bottomrule
\end{longtable}

Blaha's paper does not list separate $m\ne n$ matrix elements.  Their
Hermitian polarization is required only when the one-particle density matrices
are off-diagonal, as in DFT+eDMFT.

\section{References and source provenance}

\sloppy
\begin{enumerate}
\item K. Belbase, A. Tr\"oster, and P. Blaha, ``Stress tensor in the linearized
      augmented plane wave method,'' \emph{Phys. Rev. B} \textbf{104}, 174113
      (2021).
\item K. Haule, \emph{Forces in the LAPW implementation of DFT+eDMFT},
      internal derivation in \filepath{src/doc/Forces_LAPW.pdf}.
\item Current implementation inspected in
      \filepath{eDMFT_new/src/dmft2/cmp_dmft_stress.f90},
      \filepath{eDMFT_new/src/dmft2/cmp_dmft_weights.f90},
      \filepath{eDMFT_new/src/dmft2/l2main.f90}, and
      \filepath{eDMFT_new/src/dmft2/l2mainSubs.f90}.
\item Current audit summary:
      \filepath{eDMFT_new/src/doc/edmft_stress_audit_state_20261105.md}.
\end{enumerate}
\fussy

\vfill
\begin{center}
\small\color{bluegray}
End of technical working note.  No production source code was modified in its preparation.
\end{center}

\end{document}
